\subsection*{参考文献}

{[}1{]} W. Osgood, Non-uniform convergence and the integration of series term by term, Amer. Journ. of Math., 19 (1897), pp.~155-190.

{[}2{]} R. BAIRE, Sur les fonctions de variables réelles, Ann. di Mat. (3), 3 (1899), p.~i.

{[}3{]} H. LEBESGUE, Sur le problème des Dirichlet, Rend. Circ. mat. di Palermo, 24 (1907), pp.~371-402.

{[}4{]} H. TIETZE, Über Funktionen, die auf einer abgeschlossenen Menge stetig sind, J. de Crelle, 145 (1915), pp.~9-14.

{[}5{]} M. SOUSLIN, Sur une définition des ensembles mesurables B sans nombres transfinis, C. R. Acad. Sci., 164 (1917), pp.~88-91.

{[}6{]} P. URYSOHN, Über die Mächtigkeit der zusammenhängenden Mengen, Math. Ann., 94 (1925), p.~262.

{[}7{]} N. Lusin, Leçons sur les ensembles analytiques et leurs applications, Paris (Gauthier-Villars), 1930.

{[}8{]} K. KURATOWSKI, Topologie I, 2nd edition, Warszawa-Vrocaw, 1948.

{[}9{]} J. DIEUDONNÉ, Une généralisation des espaces compacts, Journ. de Math. (9), 23 (1944), pp.~65-76.

{[}10{]} A. H. STONE, Paracompactness and product spaces, Bull. Amer. Math. Soc., 54 (1948), pp.~977-982.

{[}11{]} G. CHOQUET, Theory of capacities, Ann. Inst. Fourier, 5 (1953-1954) pp.~131-295.

\ResumeMainChapters
